% Source-derived chapter 04: Very stable Higgs bundles
% Original source: very-stable-higgs-bundles-equivariant-multiplicity-mirror-symmetry
\chapter{Very stable Higgs bundles}
\label{very-stable-higgs-bundles-equivariant-multiplicity-mirror-symmetry-chapter-04}
\source{very-stable-higgs-bundles-equivariant-multiplicity-mirror-symmetry}

\begin{remark}[Source section]
\label{very-stable-higgs-bundles-equivariant-multiplicity-mirror-symmetry-chapter-04-source}
\source{very-stable-higgs-bundles-equivariant-multiplicity-mirror-symmetry}
\source{very-stable-higgs-bundles-equivariant-multiplicity-mirror-symmetry}
This chapter reproduces the corresponding section of the retrieved
LaTeX source, including its definitions, statements, examples, and
proofs. It has not yet been translated into Lean.
\notready
\end{remark}

% The source section title was: Very stable Higgs bundles
\label{verystable}
\source{very-stable-higgs-bundles-equivariant-multiplicity-mirror-symmetry}


\subsection{Definition and basic properties} 
Drinfeld and Laumon in \cite[Definition 3.4]{laumon} call a vector bundle $E$ on $C$ {\em very stable} if the only nilpotent Higgs field $\Phi
\in H^0(\End(E)\otimes K)$ is the trivial one.  \cite[Proposition 3.5]{laumon} proves that a rank $n$ very stable bundle is stable, and that very stable bundles form an open dense subset of the moduli space of stable  bundles. 

\begin{definition}
\source{very-stable-higgs-bundles-equivariant-multiplicity-mirror-symmetry}
\notready Let $\calE\in \M^{s \T}$ be a stable $\T$-invariant Higgs bundle. We call $\calE$ a {\em very stable Higgs bundle} if and only if  the only nilpotent Higgs bundle in $W_\calE^+$ is $\calE$.  On the other hand we call $\calE$ {\em wobbly}, when it is not very stable. 
\end{definition}

\begin{remark}
\source{very-stable-higgs-bundles-equivariant-multiplicity-mirror-symmetry}
\notready A vector  bundle which is very stable is stable, as shown in \cite[Proposition 3.5]{laumon}. For Higgs bundles we have to impose it as the definition uses upward flows in the semistable moduli space $\M$. We could also define very semistable bundles for strictly semistable bundles in the same way as above. These could also be interesting to study. 
	\end{remark}

In this language \cite[Proposition 3.5]{laumon} implies \begin{theorem} \label{laumon} There exist very stable type $(n)$ Higgs bundles $\calE\in \M^{s\T}$. In fact, they form an open dense subset of $\calN\in \pi_0(\M^{s\T})$. 
\source{very-stable-higgs-bundles-equivariant-multiplicity-mirror-symmetry}
\end{theorem}
The following was proved for very stable bundles $E$ (in our context Higgs bundles with zero Higgs field) in  \cite[Theorem 1.1]{pauly-peon-nieto}.

\begin{lemma}
\source{very-stable-higgs-bundles-equivariant-multiplicity-mirror-symmetry}
\notready\label{transversal} The stable Higgs bundle $\calE\in \M^{s\T}$ is very stable if and only if the upward flow $W^+_\calE\subset\M$ is closed.
\source{very-stable-higgs-bundles-equivariant-multiplicity-mirror-symmetry}
\end{lemma}

\begin{proof}This is Proposition~\ref{upwardclose} applied to the semi-projective variety $\M$.  
\end{proof}

\begin{remark}
\source{very-stable-higgs-bundles-equivariant-multiplicity-mirror-symmetry}
\notready \label{hsvs}
\source{very-stable-higgs-bundles-equivariant-multiplicity-mirror-symmetry}
	Recall from Remark~\ref{hitchinsection} that the upward flow    $W^+_{\calE_{L,0}}\subset\M$ of the uniformising Higgs bundle $\calE_{L,0}$  is  a Hitchin section $h:W^+_{\calE_{L,0}} \cong s(\calA)$. Thus $W^+_{\calE_{L,0}}\cap h^{-1}(0)=\{\calE_{L,0}\}$. Therefore  the  Higgs bundle $\calE_{L,0}$ is very stable. In particular, by the lemma above,  $W^+_{\calE_{{L,0}}}\subset\M$ is closed. This also follows from the fact that it is a section of the Hitchin map, which  was already pointed out in \cite[p. 454]{hitchin-lie}. 
\end{remark}

\begin{lemma}
\source{very-stable-higgs-bundles-equivariant-multiplicity-mirror-symmetry}
\notready \label{locfree} When $\calE\in \M^{\T}$ is very stable then $h:W^+_\calE\to \calA$ is finite, flat, surjective and generically \'etale. In particular, when $\calE\in \M^{\T}$ is very stable then $h_*(\calO_{W^+_\calE})$ is locally free.
\source{very-stable-higgs-bundles-equivariant-multiplicity-mirror-symmetry}
\end{lemma}
\begin{proof} As $W^+_\calE\subset \M$ is closed by Lemma~\ref{transversal} the restriction of the Hitchin map $h:W^+_\calE\to \calA$ is proper. Furthermore since $h:W^+_\calE\to \calA$ is a map between  affine varieties (cf. Proposition~\ref{bbprop}), it is quasi-finite. As a proper and quasi-finite map it is  finite (c.f \cite[Theorem 8.11.1]{ega}). Finally, it is a map between smooth equidimensional varieties and so by miracle flatness (c.f. \cite[Corollary 23.1]{matsumura}) $h$ is flat. Thus $h$ is  finite and flat  thus $h_*(\calO_{W^+_\calE})$ is locally free.
	
	Finally, as $h:W^+_\calE\to \calA$ is finite and between equidimensional varieties, it is surjective, and thus generically smooth, and so generically \'etale. 
\end{proof}

\subsection{ Hecke transformations}
\label{heckeverystable}
\source{very-stable-higgs-bundles-equivariant-multiplicity-mirror-symmetry}

For a point $c\in C$ recall the exact sequence of sheaves \beq \label{ses}0\to \calO_C(-c)\stackrel{s_c}{\to} \calO_C\to \calO_c\to 0.\eeq  Let $E\to C$ be a vector bundle. By abuse of notation we will also denote by $E$ its locally-free sheaf of sections.   Tensoring \eqref{ses} with $E$ gives \beq\label{tensored} 0\to E(-c){\to} E\stackrel{f}{\to} (E_c)_c \to 0.\eeq Here $E_c$ denotes the fibre of the vector bundle $E$ at $c$ and for any vector space $U$ we denote $U_c:=U\otimes \calO_c.$ In particular, $(E_c)_c=E\otimes \calO_c$. If now $V<E_c$ then $V_c\subset (E_c)_c$ is a subsheaf and if $W:=E_c/V$ then $W_c$ is a quotient sheaf of $(E_c)_c$. Denoting by  $\pi_V:(E_c)_c\to W_c$ the quotient map and  $f_V=f \circ \pi_V$  we make the following 
\source{very-stable-higgs-bundles-equivariant-multiplicity-mirror-symmetry}
\begin{definition}
\source{very-stable-higgs-bundles-equivariant-multiplicity-mirror-symmetry}
\notready Let $E$ be a vector bundle on the curve $C$.  Let $V<E_c$ be a subspace at a point $c\in C$. The {\em Hecke transform} of $E$ at $V$ is defined to be $E_V:=\ker f_V$. \end{definition} Equivalently, the following   short exact sequence of sheaves defines the Hecke transformed vector bundle: $$0\to E_V{\to} E\stackrel{f_V}{\to} W_c \to 0.$$ As a subsheaf of a torsion free sheaf, $E_V$ is torsion free and so a locally free sheaf, i.e. the sheaf of sections of a vector bundle which we also denote by $E_V$ or sometimes $E^\prime=E_V$. 

\begin{remark}
\source{very-stable-higgs-bundles-equivariant-multiplicity-mirror-symmetry}
\notready \label{remarkV}By construction $E(-c)\subset E_V$ is a subsheaf and fits into a short exact sequence $$0\to E(-c){\to} E^\prime{\to} V_c \to 0.$$ This way we get a surjective map $E^\prime_c\to V$. Denote by $V^\prime < E^\prime_c$ its kernel. 
\source{very-stable-higgs-bundles-equivariant-multiplicity-mirror-symmetry}
	Then we see that $E^\prime_{V^\prime}=E(-c)$. 
\end{remark}

\begin{example}
\source{very-stable-higgs-bundles-equivariant-multiplicity-mirror-symmetry}
\notready\label{heckebundle} When $V=E_c$ then the exact sequence \eqref{tensored} shows that $E_V=E(-c)$. More generally, if a vector bundle on a curve $E\oplus F$ is a direct sum of two subbundles  and we  choose the subspace $E_c<(E\oplus F)_c$ then the Hecke transfrom of $E\oplus F$ at $E_c$ satisfies $$(E\oplus F)_{E_c}\cong E(-c)
\source{very-stable-higgs-bundles-equivariant-multiplicity-mirror-symmetry}
\oplus F.$$ This can be seen by adding the trivial sequence $0\rightarrow F\stackrel{\cong}{\rightarrow}F\to 0 \to 0$ to \eqref{tensored}.	\end{example}

Let now $(E,\Phi)$ be a Higgs bundle and let $V<E_c$ be a $\Phi_c$-invariant subspace. Consider the quotient $W:=E_c/V$. Denote by $\overline{\Phi}_c$ the morphism $W_c\to W_c\otimes K$ induced by $\Phi$.   This data induces a unique Higgs bundle $(E^\prime,\Phi^\prime)$ making the following diagram commutative.
\beq \label{HeckeHiggs} \begin{array}{ccccccc} 0 \to & E^\prime& {\to} &E& \stackrel{f_V}{\to} & W_c & \to 0 
\source{very-stable-higgs-bundles-equivariant-multiplicity-mirror-symmetry}
	\\ & \Phi^\prime \downarrow &&\Phi \downarrow&& \overline{\Phi}_c\downarrow & \\ 0 \to & E^\prime K&\to &EK&\to& W_c K&\to 0 \end{array}.\eeq 

\begin{definition}
\source{very-stable-higgs-bundles-equivariant-multiplicity-mirror-symmetry}
\notready Let $(E,\Phi)$ be a Higgs bundle on the curve $C$ and $V< E_c$ a $\Phi_c$-invariant subspace at $c\in C$. The {\em Hecke transform} of $(E,\Phi)$ at $V< E_c$ is the unique Higgs bundle $(E^\prime,\Phi^\prime)$ making the diagram \eqref{HeckeHiggs} commutative. Sometimes we will denote it $\calH_V(E,\Phi)=(E^\prime,V^\prime)$ or by $(E_V,\Phi_V)=(E^\prime,\Phi^\prime)$. 
\end{definition}
% Clearly $\deg(E^\prime)=\deg(E)-n+k$.  

\begin{remark}
\source{very-stable-higgs-bundles-equivariant-multiplicity-mirror-symmetry}
\notready \label{V'Higgs}As  in Remark~\ref{remarkV}, the Hecke transformation will produce a subspace $V^\prime<E^\prime_c$ such that the Hecke transform of $E^\prime$ at $V^\prime$ is $E(-c)$. We note that the Hecke transformed Higgs field $\Phi^{\prime\prime}:E(-c)\to E(-c)K$ as  a section of $\End(E(-c))K=\End(E)K$ agrees with $\Phi\in \End(E)K$ away from $c$, and so agrees everywhere. Thus the Hecke transform of $(E^\prime,\Phi^\prime)$ at $V^\prime$ is the Higgs bundle $(E(-c),\Phi)$. 
\source{very-stable-higgs-bundles-equivariant-multiplicity-mirror-symmetry}
\end{remark}

\begin{remark}
\source{very-stable-higgs-bundles-equivariant-multiplicity-mirror-symmetry}
\notready\label{alonghitchin} We also note that as $\Phi^\prime$ and $\Phi$ agree away from $c$ so  their characteristic polynomial will also agree away from $c$, therefore $h(\Phi^\prime)=h(\Phi)\in \oplus_{i=1}^n H^0(C;K^i)$. Therefore the Hecke transformation is along the fibers of the Hitchin fibration. In the case of a point $a\in {\mathcal A}$ for which  the spectral curve $C_a$ is smooth, the fibre is isomorphic to the Jacobian of $C_a$. At a generic point $c\in C$ the subspace $V<E_c$ is a sum of  eigenspaces of $\Phi_c$. Since $C_a$ is the curve of eigenvalues this is an effective  divisor $D$. In the notation of Section~\ref{intersection} the Hecke transform takes the line bundle $U$ on $C_a$ corresponding to $(E,\Phi)$  to the line bundle $U(-D)$ corresponding to $(E^\prime,\Phi^\prime)$. 
\source{very-stable-higgs-bundles-equivariant-multiplicity-mirror-symmetry}
\end{remark}

\begin{example}
\source{very-stable-higgs-bundles-equivariant-multiplicity-mirror-symmetry}
\notready\label{invarianthecke} Assume that the Higgs bunde $(E,\Phi)$ is such that $E=E_1\oplus \dots \oplus E_k$, a direct sum of subbundles, and $\Phi(E_i)\subset E_{i+1}K$. Then for any $1\leq j \leq k-1$ we have  $V:=(E_{j+1}\oplus \dots \oplus E_k)_c<E_c$ is $\Phi_c$ invariant. Then we can determine the Hecke transform $(E^\prime,\Phi^\prime)$ at $V$ by definition as $E^\prime=E_1^\prime \oplus \dots \oplus E^\prime_k$ such that $E_i^\prime=E_i(-c)$ unless $i>j$ when $E_i^\prime=E_i$ from Example~\ref{heckebundle}.  The modified Higgs field then satisfies $$\Phi^\prime|_{E^\prime_i}=\Phi|_{E_i}\in H^0(C;E_i^{\prime*}E^\prime_{i+1}K)\cong H^0(C;E_i^*E_{i+1}K)$$ for all $i$ except $i=j$ when $$\Phi^\prime|_{E_{j}(-c)}=\Phi|_{E_j}s_c\in H^0(C;E_i^*(c)E_{i+1}K).$$ In particular, $(\Phi^\prime|_{E_{j}(-c)})$ vanishes at $c$. 
\source{very-stable-higgs-bundles-equivariant-multiplicity-mirror-symmetry}
Furthermore $V^\prime$ of Remark~\ref{V'Higgs} will be $V^\prime=E^\prime_1\oplus \dots \oplus E^\prime_{j}|_c<E^\prime_c$, so that the $\calH_{V^\prime}(E^\prime,\Phi^\prime)=(E(-c),\Phi)$ as expected.

Doing the Hecke transformation backwards we get that if the Higgs bundle $(E,\Phi)$ has a compatible decomposition $E=E_1\oplus\dots\oplus E_k$ such that $\Phi(E_i)\subset E_{i+1}K$ and $\Phi_c$ is zero on $(E_j)_c$, then $V=(E_1\oplus \dots \oplus E_{j})|_c<E_c$ is $\Phi_c$-invariant.  Then the Hecke transform
$\calH_{V}(E,\Phi)$ has the chain form\footnote{Where arrows represent morphisms twisted by $K$.}
$$E_1\stackrel{\Phi_1}{\longrightarrow}{E_2}\stackrel{\Phi_2}{\longrightarrow}\cdots \stackrel{\Phi_{j-2}}{\longrightarrow}E_{j-1} \stackrel{\Phi_{j-1}/s_c}{\longrightarrow} E_j(-c)\stackrel{\Phi_{j}}{\longrightarrow}\cdots \stackrel{\Phi_{k}}{\longrightarrow}{E_k(-c)},$$ where $\Phi_i:=\Phi|_{E_i}:E_i\to E_{i+1}K$. 

Note that in both cases above the Higgs bundles are $\T$-fixed and so are the Hecke transformed Higgs bundles (when all of them are stable). However if the $\Phi_c$-invariant subspace is not a direct sum of the $E_i|_c$'s then  the Hecke transform will not be $\T$-fixed. Such examples will play a crucial rule in the only if part of the proof of Theorem~\ref{mainverystable}. 

	\end{example}

%We can recover $(E,\Phi)$ from $(E^\prime)$ and the $n-k$ dimensional subspace $\ker{\iota_x}\subset E^\prime_x$ as follows.  Its annihilator in the dual space is denoted by $V^\prime:=ann(\ker(\iota_x)) \subset E^{\prime *}_x$, which has dimension $k$. Finally denote the quotient $W^\prime:=E^{\prime *}_x/V^\prime$, an $n-k$ dimensional vector space.  Take $E^*\stackrel{\iota^*}{\to} E^{\prime *}$. It is an isomorphism away from $x$ and thus is an injective map of sheaves. It fits into the following commutative diagram:   
%\bes \begin{array}{ccccccc} 0 \to & E^{ *} & \stackrel{\iota^*}{\to} &E^{\prime *}& \stackrel{f_V^\prime}{\to} & W^\prime_x & \to 0 
%\\ & \Phi^{*} \downarrow &&\Phi^{\prime*} \downarrow&& \overline{\Phi}^{\prime *}_x\downarrow & \\ 0 \to & E^{*} K&\to &E^{\prime *}K&\to& W^\prime_x K&\to 0 \end{array}.\ees

%Thus $(E^*,\Phi^*)$ is the Hecke transform of $(E^{\prime*},\Phi^{\prime*})$ at $V^\prime< E^{\prime*}_{x}$. 

\subsubsection{The Hecke transform of a full compatible filtration}

Here we investigate what happens to a compatible full filtration on a Higgs bundle,  when we perform a Hecke transformation at a point $c\in C$ where the Higgs field on the fibre of the graded Higgs bundle at $c$ is regular nilpotent. This is equivalent with $b(c)\neq 0$ using the notation \eqref{be} below. 

Let a rank $n$ Higgs bundle $(E,\Phi)$ carry a full filtration \beq\label{fullfiltration}0=E_0\subset E_1\subset\dots\subset E_n=E\eeq by $\rank(E_i)=i$ subbundles, such that \beq\label{compatible}\Phi(E_i)\subset E_{i+1}\otimes K.\eeq Then we have the induced maps \beq\label{bi}b_i=:E_i/E_{i-1}\to E_{i+1}/E_i\otimes K\eeq between line bundles, and the composition \beq \label{be}b=b_{n-1}\circ \dots \circ b_1:E_1\to E/{E_{n-1}}\otimes K^{n-1}.\eeq Suppose that $b$ does not vanish at $c\in C$. This means that for all $i$ the map  $\Phi_c: E_{i}|_c\to E_{i+1}|_c\otimes K|_c$ does not preserve $E_i|_c$.  To perform the Hecke transformation we choose a $\Phi_c$-invariant subspace $V<E_c$ of $\dim(V)=k< n$.  We need the following linear algebra
\source{very-stable-higgs-bundles-equivariant-multiplicity-mirror-symmetry}
\begin{lemma}
\source{very-stable-higgs-bundles-equivariant-multiplicity-mirror-symmetry}
\notready\label{linear} Let $0=U_0<U_1<\dots U_n<U$ be a full flag of subspaces of an $n$-dimensional complex vector space $U$. Furthermore let $A$ be a linear transformation $A:U\to U$ such that $A(U_i)<U_{i+1}$ but $U_i$ is not $A$-invariant. If $V< U$ of $\dim(V)=k< n$ is $A$-invariant, then 
\source{very-stable-higgs-bundles-equivariant-multiplicity-mirror-symmetry}
	$\dim (V\cap U_i)= \max(0,k+i-n)$.
\end{lemma}
\begin{proof} By assumption $U_i+A(U_i)= U_{i+1}$. Thus $V+U_i=V+U_{i+1}$ if and only if $V+U_i$ is invariant under $A$. On the other hand $V+U_i=V+U_{i+1}$ in turn implies $V+U_i=V+U_n=U$. 
	Thus we see that the chain of subspaces
	$$V=V+U_0\leq V+U_1\leq \dots \leq V+U_{n-1}\leq V+U_n=U$$ must be strictly increasing until it reaches the whole ambient space $U$. Therefore $\dim(V+U_i)=\min(k+i,n)$, which implies the result.  
\end{proof}

The $\Phi_c$-invariant subspace $V<E_c$ induces a map of sheaves $f_V:E\to W_c$, where $W:=E_c/V$.   This induces the Hecke transformed Higgs bundle $(E^\prime,\Phi^\prime)$:
\bes \begin{array}{ccccccc} 0 \to & E^\prime & \to &E& \stackrel{f_V}{\to} & W_c & \to 0 
	\\ & \Phi^\prime \downarrow &&\Phi \downarrow&& \overline{\Phi}_c\downarrow & \\ 0 \to & E^\prime K&\to &EK&\to& W_c K&\to 0 \end{array}.\ees 
Each subbundle $E_i$ will also induce a Hecke transformed line bundle with respect to the subspace $V_i:=E_i|_c\cap V< E_i|_c$. Denoting $W^i:=E_i|_c/V_i$ we get a new vector bundle $E_i^\prime$ from the kernel of the map $f_{V_i}:E_i\to W^i_c$. In other words we have the top row of the following commutative diagram \beq \label{commutative} \begin{array}{ccccccc} 0\to & E^\prime_i&\to &E_i& \stackrel{f_{V_i}}{\to} &W^i_c&\to 0 \\ &\rotatebox[origin=c]{-90}{$\hookrightarrow$} &&\rotatebox[origin=c]{-90}{$\hookrightarrow$} &&\rotatebox[origin=c]{-90}{$\hookrightarrow$}& \\   0\to & E_{i+1}^\prime&\to &E_{i+1}& \stackrel{f_{V_{i+1}}}{\to} &W^{i+1}_c&\to 0 \\ &\rotatebox[origin=c]{-90}{$\twoheadrightarrow$} &&\rotatebox[origin=c]{-90}{$\twoheadrightarrow$} &&\rotatebox[origin=c]{-90}{$\twoheadrightarrow$}& \\   0\to & E_{i+1}^\prime/E^\prime_i&\to &E_{i+1}/E_i& \stackrel{\overline{f}_{V_{i+1}}}{\to} &(W^{i+1}/W^i)_c&\to 0
\source{very-stable-higgs-bundles-equivariant-multiplicity-mirror-symmetry}
\end{array}\eeq Here the embedding  $E_i^\prime\hookrightarrow E_{i+1}^\prime$ is uniquely induced from the rest of the first two top rows of the diagram. The quotient maps connecting the second and third rows of the diagram are the ones making the columns  into short exact sequences. The last row, and in particular the map  $\overline{f}_{V_i}$,  is then  uniquely induced from the rest of the diagram to make it commutative. 

From Lemma~\ref{linear} we have  \beq\label{dimwi}\dim(W^{i+1}/W^i) =\left\{
\source{very-stable-higgs-bundles-equivariant-multiplicity-mirror-symmetry}
\begin{array}{cc}1 & \mbox{ for } i< n-k \\ 0 & \mbox{ for } i\geq n-k \end{array}
\right.\eeq
Thus from the bottom row of \eqref{commutative} we obtain \beq \label{'quotients}E_{i+1}^\prime/E^\prime_i\cong\left\{
\source{very-stable-higgs-bundles-equivariant-multiplicity-mirror-symmetry}
\begin{array}{cc}(E_{i+1}/E_i)(-c) & \mbox{ for } i< n-k \\ (E_{i+1}/E_i) & \mbox{ for } i\geq n-k \end{array}
\right.\eeq In particular, $E_{i+1}^\prime/E^\prime_i$ is a line bundle, and so $E^\prime_i\subset E^\prime_{i+1}$ is a subbundle. Thus we get a full filtration of $E^\prime$ by  subbundles: \beq \label{'filtration} E^\prime_0=0\subset E^\prime_1\subset\dots\subset E^\prime_n=E
\source{very-stable-higgs-bundles-equivariant-multiplicity-mirror-symmetry}
\eeq

By construction it is compatible with $\Phi^\prime$ i.e.  \beq \label{'compatible}\Phi^\prime(E^\prime_i)\subset E^\prime_{i+1}K.\eeq Finally we will determine the maps $b^\prime_i:E^\prime_i/E^\prime_{i-1}\to E^\prime_{i+1}/E^\prime_i\otimes K$ induced from $\Phi^\prime$. We have the commutative diagram
\source{very-stable-higgs-bundles-equivariant-multiplicity-mirror-symmetry}

\bes \begin{array}{ccccccc}   0\to & E_{i}^\prime/E^\prime_{i-1}&\to &E_{i}/E_{i-1}& \stackrel{\overline{f}_{V_{i}}}{\to} &(W^{i}/W^{i-1})_c&\to 0 \\ &b_i^\prime \downarrow &&b_i \downarrow &&\overline{b}_i\downarrow & \\   0\to & (E_{i+1}^\prime/E^\prime_i)\otimes K&\to &(E_{i+1}/E_i)\otimes K& \stackrel{\overline{f}_{V_{i+1}}}{\to} &(W^{i+1}/W^i)_c\otimes K&\to 0
\end{array}.\ees Here $\overline{b}_i$ is induced from $\Phi_c$. When $i>n-k$ then from \eqref{dimwi} we see that $\overline{b}_i$ maps between zero-dimensional spaces, which shows that \beq \label{bibip}b_i^\prime=b_i\in H^0(C;(E_{i}^\prime/E^\prime_{i-1})^*(E_{i+1}^\prime/E^\prime_i)K)=H^0(C;(E_{i}/E_{i-1})^*(E_{i+1}/E_i)K).
\source{very-stable-higgs-bundles-equivariant-multiplicity-mirror-symmetry}
\eeq
When $i<n-k$ then $\overline{b}_i$ is induced from an isomorphism of $1$-dimensional vector spaces, and we get again the agreement \eqref{bibip}. 

Furthermore, when $i=n-k$ then $W^{i}/W^{i-1}$ is one-dimensional and $W^{i+1}/W^i=0$. This implies that $b_i^\prime=b_i s_c$ where $s_c\in H^0(C;\calO(c))\cong H^0(C;\Hom(\calO(-c), \calO))$ is the defining \eqref{ses} section vanishing at $c$. 

Finally by tracing the kernel of the induced maps on the fibres at $c$ in the diagram \eqref{commutative} and using \eqref{dimwi} we see that $V^\prime=(E^\prime_{n-k})_c<E^\prime_c$. Thus $\calH_{V^\prime}(E^\prime,\Phi^\prime)=(E(-c),\Phi )$.

We proved the following

\begin{proposition}
\source{very-stable-higgs-bundles-equivariant-multiplicity-mirror-symmetry}
\notready \label{modifiedfilt}Let a Higgs bundle $(E,\Phi)$ carry a full \eqref{fullfiltration} and compatible \eqref{compatible} filtration, such that $b$ \eqref{be} does not vanish at $c$. Let $V<E_c$ be a  $k$-dimensional and $\Phi_c$-invariant subspace. Then the Hecke transformed Higgs bundle $\calH_V(E,\Phi)=(E^\prime,\Phi^\prime)$ carries an induced full filtration \eqref{'filtration} which is compatible with $\Phi^\prime$ \eqref{'compatible}. Furthermore we have \eqref{'quotients} and $b^\prime_i=b_i$ unless $i=n-k$ when $b^\prime_{n-k}=b_{n-k}s_c$ acquires a simple zero at $c$. Finally, the subspace of Remark~\ref{V'Higgs} is  $V^\prime=(E^\prime_{n-k})_c<E^\prime_c$. Thus $\calH_{V^\prime}(E^\prime,\Phi^\prime)=(E(-c),\Phi )$. 
\source{very-stable-higgs-bundles-equivariant-multiplicity-mirror-symmetry}
\end{proposition}



\subsubsection{Very stable Higgs bundles via Hecke transformations}



For $\delta=(\delta_0,\delta_1,\dots,\delta_{n-1})\in J_\ell(C)\times C^{[m_1]}\times \dots \times C^{[m_{n-1}]}$ let $\calE_{\d}\in \M^{s\T}$ be a type $(1,...,1)$ fixed point of the $\T$-action constructed in Example~\ref{ex111}.

\begin{theorem}
\source{very-stable-higgs-bundles-equivariant-multiplicity-mirror-symmetry}
\notready \label{mainverystable} Let $\calE_{\d} \in \M^{s\T}$ be a type $(1,...,1)$ fixed point of the $\T$-action. Then $\calE_{\d}$ is very stable if and only if $b=b_1\circ\dots \circ b_{n-1}\in H^0(C;L_0^*L_{n-1}K)$ has no repeated zero, i.e. if and only if the effective divisor $\div(b)=\delta_1+\dots+\delta_{n-1}$ is reduced. 
\source{very-stable-higgs-bundles-equivariant-multiplicity-mirror-symmetry}
\end{theorem}
\begin{proof} First we prove the if part. The proof is by induction on $\deg(b)$. When $b$ has no zeroes, then $\calE_{\d}\cong \calE_{L_0,0}$ are the uniformising Higgs bundles from Remark~\ref{hitchinsection}. Thus $W^+_{\calE_{\d}}$ is a Hitchin section. In particular $\calE_{\d}$ is its intersection with the nilpotent cone $h^{-1}(0)$ (c.f. Remark~\ref{hsvs}). Thus $\calE_{\d}$ is very stable when $\deg(b)=0$. 
	
	Let $\calE_{\d} \in \M^{s\T}$ be such that $b$ has no repeated zero, but at least one zero say $c\in C$ with $b(c)=0$. Assume that it is $b_k(c)=0$. Then $V:=(L_0\oplus \dots \oplus L_{k-1})_c<E_c$ is $\Phi_c$ invariant. We see from Example~\ref{invarianthecke} that $\calH_{V}(\calE_{\d})=\calE_{\d^\prime}$ where $\d^\prime_i=\d_i$ i.e. $b^\prime_i=b_i$ except if  $i=k$ when  $b^\prime_k=b_k/s_c$ i.e. $\d^\prime_k=\d_k-c$. In particular, $b^\prime$ does not vanish at $c$. We will make frequent use of the following
	\begin{lemma}
\source{very-stable-higgs-bundles-equivariant-multiplicity-mirror-symmetry}
\notready\label{stablemma} Let $\calE_{\d} \in \M^{s\T}$ such that $b_k(c)=0$. Let $V:=(L_0\oplus \dots \oplus L_{k-1})_c<E_c$. Then  $\calE_{\d^\prime}:=\calH_{V}(\calE_{\d})$ has the chain form
\source{very-stable-higgs-bundles-equivariant-multiplicity-mirror-symmetry}
$$L_0\stackrel{b_1}{\rightarrow}{L_1}\stackrel{b_2}{\rightarrow}\cdots \stackrel{b_{k-1}}{\rightarrow}L_{k-1} \stackrel{b_k/s_c}{\rightarrow} L_k(-c)\stackrel{b_{k+1}}{\rightarrow}\cdots \stackrel{b_{n-1}}{\rightarrow}{L_{n-1}(-c)}$$ and it is a stable Higgs bundle.
	\end{lemma}
	\begin{proof} As by assumption $\calE_{\d}$ is stable all proper $\Phi$-invariant subbundles have slope less than the slope of $E$. As the  only $\Phi$-invariant subbundles are $L_i\oplus \dots \oplus L_{n-1}\subset E$, we have for every $0\leq i \leq n-1$ the following inequality  $$\frac{\ell_i+\dots+\ell_{n-1}}{n-i}<\frac{\ell_0+\dots+\ell_{n-1}}{n}.$$ To see that $\calE_{\d^\prime}$ is stable we check that \bes \frac{\ell^\prime_i+\dots+\ell^\prime_{n-1}}{n-i}=\frac{\ell_i+\dots+\ell_{n-1}-\min(n-i,n-k)}{n-i}&<&\frac{\ell_0+\dots+\ell_{n-1}}{n}-\frac{\min(n-i,n-k)}{n-i}\\
		< \frac{\ell^\prime_0+\dots+\ell^\prime_{n-1}+n-k}{n}-\frac{\min(n-i,n-k)}{n-i}	
		&<& \frac{\ell^\prime_0+\dots+\ell^\prime_{n-1}}{n} \ees 
		
		because $$\frac{n-k}{n}\leq \frac{\min(n-i,n-k)}{n-i}.$$
	\end{proof}
	
	Assume now that we have a nilpotent Higgs bundle $\calE=(E,\Phi)\in W^+_{\calE_{\d}}$. From Proposition~\ref{filtrationexists} let $E_1\subset\dots\subset E_n=E$ be the $\Phi$-invariant filtration  inducing $\lim_{\lambda\to 0}\lambda\cdot \calE=\calE_{\d}$. Let $V:=(E_{k})_c<E_c$, which by the assumption of $b_k(c)=0$ is $\Phi_c$-invariant. As we saw in the previous subsection $\calH_V(E,\Phi)=(E^\prime,\Phi^\prime)$ is nilpotent in the upward flow of $\calH_{V_{1k}}(\calE_{\d})=\calE_{\d^\prime}$, because $\calE_{\d^\prime}$ is stable by Lemma~\ref{stablemma}. But $\deg(b^\prime)=\deg(b)-1$ and so by induction we can assume that $\calE_{\d^\prime}$ is very stable and hence the nilpotent $(E^\prime,\Phi^\prime)$ in its upward flow should be isomorphic to it: $(E^\prime,\Phi^\prime)\cong \calE_{\d^\prime}$. Now $b^\prime(c)\neq 0$ therefore $\Phi^\prime_c$ is regular nilpotent. Hence there is a unique $\Phi^\prime_c$-invariant $(n-k)$-dimensional subspace of $E^\prime$, namely $V^\prime_{k,n-1}=(L_{k})_c\oplus \dots\oplus (L_{n-1})_c$. The Hecke transform of $(E^\prime,\Phi^\prime)\cong \calE_{\d^\prime}$ at this subspace gives $(E(-c),\Phi)\cong\calE_{\delta-c}$, where $$\delta-c:=(\delta_0-c,\delta_1,\dots,\delta_{n-1}).$$ This shows that $\calE_{\d}$ is also very stable. This proves the ``if" part of the theorem.
	
	For the only if part, when $b$ has repeated zeroes, the strategy will be to produce a {\em Hecke curve} \cite{NR} - a one-parameter family of Higgs bundles produced as Hecke transforms of a fixed Higgs bundle -   in the intersection of the nilpotent cone and the upward flow $W^+_{\calE_{\d}}$ originating at $\calE_{\d}$. 
	
	Take a multiple zero $c$ of $b$. Assume that $b_k(c)=0$. Then $V:=(L_0\oplus \dots \oplus L_{k-1})_c<E_c$ is $\Phi_c$-invariant. The Hecke transformed Higgs bundle from Example~\ref{invarianthecke} $\calH_{V}(\calE_{\d})=\calE_{\d^\prime}=(E^\prime,\Phi^\prime)$ has  $b^\prime(c)=0$ otherwise $\calH_{V^\prime}(\calE_{\d^\prime})$ would have $b(c)=0$ a simple zero from Proposition~\ref{modifiedfilt}. 
	
	Thus $b^\prime(c)=0$ and as before in Lemma~\ref{stablemma} $\calE_{\d^\prime}$ is still stable. Now however the nilpotent $\Phi^\prime_c$ is no longer regular due to $b^\prime(c)=0$. This implies that, besides $V^\prime=(L^\prime_{k}\oplus \dots \oplus L^\prime_{n-1})_c$,  
	there is more than one $\Phi^\prime_c$-invariant $n-k$-dimensional subspace in $E^\prime_c$.  By \cite[Theorem 6]{shayman} the subvariety $S_{n-k}(\Phi^\prime_c)\subset Gr_{n-k}(E^\prime_c)$  of $\Phi^\prime_c$-invariant $(n-k)$-dimensional subspaces in the Grassmannian  of $(n-k)$-planes in $E^\prime_c$ is connected. The connectedness of $S_{n-k}(\Phi^\prime_c)$ also follows from \cite[Theorem 6.2]{horrocks} as one can think of $S_{n-k}(\Phi^\prime_c)$ as the support of the fixed point scheme of the unipotent group action of $\G_a\cong \{\exp(t \Phi^\prime_c)\}$ on the Grassmannian $Gr_{n-k}(E^\prime_c)$.
	
	We will find a curve in $ S_{n-k}(\Phi^\prime_c)$, leading to a Hecke curve in $\M$ by studying a natural $\T$-action on the projective variety $S_{n-k}(\Phi^\prime_c)$. We note that $\calE_{\d^\prime}=(E^\prime, \Phi^\prime)$ is $\T$-invariant, meaning that \beq \label{isom}(E^\prime, \Phi^\prime)\cong (E^\prime, \lambda \Phi^\prime)\eeq for all $\lambda\in \C$. The isomorphism in \eqref{isom} is induced by an automorphism of $(E^\prime,\Phi^\prime)$. Let $\T$ act on $E^\prime=L^\prime_0\oplus \dots \oplus L^\prime_{n-1}$ with weight $-i$ on $L^\prime_i$. This $\T$-action induces the weight $-1$ action of $\T$ on $\Phi^\prime$, showing \eqref{isom}. Restricting this $\T$-action to $E^\prime_c$ we get an action of $\T$ on $Gr_{n-k}(E^\prime_c)$. We note that $\Phi^\prime_c(\lambda\cdot v)= \lambda\cdot \Phi^\prime_c(v)$ and thus if $V\subset E^\prime_c$ is a $\Phi_c$-invariant subspace then $\lambda\cdot V$ is also $\Phi_c$-invariant. Therefore this $\T$-action leaves $S_{n-k}(\Phi^\prime_c)$ invariant. We have $S_{n-k}(\Phi^\prime_c)^\T\subset Gr_{n-k}(E^\prime_c)^\T$.  Fixed points of the $\T$-action on $Gr_{n-k}(E^\prime_c)$ are of the form $(L_{i_1}\oplus \dots \oplus L_{i_{n-k}})_c$ for any $(n-k)$-element subset $\{i_1,\dots,i_{n-k} \}\subset \{0,\dots,n-1\}$. We  have $$V^\prime=(L^\prime_{k},\dots, L^\prime_{n-1})_c\in S_{n-k}(\Phi^\prime_c)^\T\subset Gr_{n-k}(E_c^\prime).$$ By Remark~\ref{V'Higgs} we have that $\calH_{V^\prime}(E^\prime,\Phi^\prime)=(E(-c),\Phi)=\calE_{\delta-c}$. 
\source{very-stable-higgs-bundles-equivariant-multiplicity-mirror-symmetry}

 As $S_{n-k}(\Phi^\prime_c)$ is projective, connected and contains more than one point, there is $V^\prime\neq V\in S_{n-k}(\Phi^\prime_c)$ such that $\lim_{\lambda\to 0} \lambda V=V^\prime$. Let $$V_\infty:=\lim_{\lambda\to \infty} \lambda\cdot V\in S_{n-k}(\Phi^\prime_c)^\T.$$ As the $\T$-action on $S_{n-k}(\Phi^\prime_c)$ is linear we see that  $V^\prime\neq V_\infty$. Let $V_\infty=(L_{i_1}\oplus \dots \oplus L_{i_{n-k}})_c$ for some \beq\label{subset}\{k,\dots, n-1\}\neq \{i_1,\dots,i_{n-k} \}\subset \{0,\dots,n-1\}.\eeq We thus have a map  $\P^1\to S_{n-k}(\Phi^\prime_c)$ sending $0$ to $V^\prime$, $\infty$ to $V_\infty$ and $\lambda \in \T$ to $\lambda\cdot V$. We will denote the corresponding subspaces $V_t$ for $t\in \C \cup {\infty}$. 
\source{very-stable-higgs-bundles-equivariant-multiplicity-mirror-symmetry}
	
	First we need the following 
	
	\begin{lemma}
\source{very-stable-higgs-bundles-equivariant-multiplicity-mirror-symmetry}
\notready $\calH_{V_\infty}(\calE_{\d^\prime})$ is a $\T$-fixed stable Higgs bundle. 
	\end{lemma}
	\begin{proof} Let $\calH_{V_\infty}(\calE_{\d^\prime})=(E^{\prime\prime},\Phi^{\prime\prime})$. Then $E^{\prime\prime}\cong L_0^{\prime\prime}\oplus \dots \oplus L_{{n-1}}^{\prime\prime}$ where $L_i^{\prime\prime}\cong L_i^\prime$ when $i\in\{i_1,\dots,i_{n-k}\}$ and $L_i^{\prime\prime}\cong L_i^\prime(-c)$ when $i\notin\{i_1,\dots,i_{n-k}\}$. Thus 
		\beq\label{lpp}  \ell_i^{\prime\prime}:=\deg(L_i^{\prime\prime} ) =\left\{ \begin{array}{cc} \ell_i^\prime & \mbox{ when } i\in\{i_1,\dots,i_{n-k}\} \\ \ell_i^\prime-1 & \mbox{ when } i\notin\{i_1,\dots,i_{n-k}\}  
\source{very-stable-higgs-bundles-equivariant-multiplicity-mirror-symmetry}
		\end{array} 
		\right. \eeq
		The only $\Phi^{\prime\prime}$-invariant subbundles of $E^{\prime\prime}$ are of the form $L_i^{\prime\prime}\oplus \dots \oplus L_{n-1}^{\prime\prime}$. Thus to prove stability we check \bes  \frac{\ell^{\prime\prime}_i+\dots+\ell^{\prime\prime}_{n-1}}{n-i}\leq\frac{\ell^\prime_i+\dots+\ell^\prime_{n-1}-\max(0,k-i)}{n-i}&=&\frac{\ell_i+\dots+\ell_{n-1}-(n-i)}{n-i}  \\
		&<&  \frac{\ell_0+\dots+\ell_{n-1}-n}{n}=
		\frac{\ell^{\prime\prime}_0+\dots+\ell^{\prime\prime}_{n-1}}{n}
		\ees
		
	\end{proof}
	
	Let $V\in S_{n-k}(E_c)$ and denote $V_\lambda:=\lambda\cdot V$. Denote by $f_\lambda:E^\prime\to E^\prime$ the automorphism induced by our $\T$-action on $E^\prime$ as in \eqref{fixed}. Finally identify  $(K)_c\cong \C$.   We now have the commutative diagram: 
	
	
	$$\begin{tikzcd}[column sep={between origins,5em}]
	& E^{\prime}_VK \ar{rr}
	& & E^\prime K \ar{rr} & & (E_c/V)_c \\
	E^{\prime}_V \ar{ur}{\lambda \Phi^{\prime}_V} \ar[crossing over]{rr} \ar{dd}[swap,near end]{f^{\prime\prime}_\lambda} 
	& & E^{\prime} \ar{ur}{\lambda \Phi^\prime}\ar{rr} & & (E_c/V)_c \ar{ur}{\lambda \overline{\Phi}^\prime_c}& \\
	& E^{\prime}_{V_\lambda}K \ar[leftarrow]{uu}[near start]{f^{\prime\prime}_\lambda}  \ar{rr} \ar[leftarrow, swap]{dl}{  \Phi^{\prime}_{V_\lambda} }
	& & E^\prime K \ar[leftarrow, swap]{dl}{ \Phi^\prime} \ar{rr} \ar[leftarrow, near start]{uu}{f^\prime_\lambda}& & (E_c/V_\lambda)_c \ar[leftarrow,near start]{uu}{\overline{f}^\prime_{\lambda c}}\\
	E^\prime_{V_\lambda} \ar{rr} & & E^\prime \ar[leftarrow, near start]{uu}{f^\prime_\lambda} \ar{rr} & & (E_c/V_\lambda)_c \ar[leftarrow, near start]{uu}{\overline{f}^\prime_{\lambda c}} \ar{ur}{ {\overline{\Phi}^\prime_c}} &
	\end{tikzcd}.$$
	
	The map $f_\lambda^{\prime\prime}:E_V^\prime\to E^\prime_{V_\lambda}$ is defined uniquely so as to make the diagram commutative.  The diagram shows that $f_\lambda^{\prime\prime}$ induces $(E^\prime_{V_\lambda},\Phi^\prime_{V_\lambda})\cong (E^\prime_V,\lambda\cdot \Phi^\prime_V)$. 
	
	Therefore   $\calH_{ V_t}(\calE_{\d})$ is stable for all $t$ since stability is an open condition for Higgs bundles (c.f. \cite[Lemma 3.7]{simpson1} or \cite[Proposition 3.1]{nitsure}). Finally, we note that $  \calH_{ V_0}(\calE_{\d})\cong \calE_{\d-c} \not\cong \calH_{ V_\infty}(\calE_{\d})$ because there exists an $i$ such that $\deg(L_i^{\prime\prime})\neq \deg(L_i(-c))$ because of \eqref{subset} and \eqref{lpp}. Thus we found a $\T$-invariant family of nilpotent stable Higgs bundles in $\M^s$ parametrized by $\bP^1$ connecting  $\calE_{\d-c}$ and $\calH_{ V_\infty}(\calE_{\d})$. Thus $W^+_{\calE_{\d-c}}\cap h^{-1}(0)$ is not trivial. Therefore $\calE_{\d-c}$ and so $\calE_{\d}$ is not very stable. The  theorem follows. 
	
\end{proof}

\begin{corollary}
\source{very-stable-higgs-bundles-equivariant-multiplicity-mirror-symmetry}
\notready \label{type111} There is a very stable Higgs bundle in every type $(1,\dots,1)$ component $F_{\mm}\in \pi_0(\M^{s\T})$ from \eqref{type111fix}. In fact, they form a dense open subset. 
\source{very-stable-higgs-bundles-equivariant-multiplicity-mirror-symmetry}
\end{corollary}

\begin{proof} From Theorem~\ref{mainverystable} the very stable locus in $F_{\mm}$ are Higgs bundles $\calE_{\d}$ where $b$ has no repeated zero. Thus they form the complement of some divisor. The result follows. 
\end{proof}

\subsubsection{Example in rank $2$}

Let $L_0$ and $L_1$ be line bundles on $C$ and $E=L_0\oplus L_1$ together with a lower triangular Higgs field: $\Phi=\left(\begin{array}{cc}0&0\\b&0\end{array} \right)$ with $b=b_1\in H^0(C;L_0^*L_1K)$. Assume that $\deg(b)<2g-2$ so that $(E,\Phi)$ is stable. Suppose further that $b=s_c^2 b^{\prime\prime}$ has a double zero at $c\in C$ where $b^{\prime\prime}\in H^0(C;L^*_0L_1K(-2c)).$ Let $V:=(L_0)_c< E_c$. It is $\Phi_c$-invariant as $\Phi_c=0$. We have $$E^\prime=L_0^\prime\oplus L_1^\prime=L_0\oplus L_1(-c)$$ and $V^\prime=V_0= (L_1^\prime)_c$. As $b^\prime=s_c b^{\prime\prime}$ still vanishes at $c$, $\Phi^\prime_c=0$. Thus $S_{1}(\Phi^\prime_c)=\P(E_c)$ the whole $\P^1$. Let $V_\infty=(L_0)_c<E^\prime_c$, which is $\Phi^\prime_c$-invariant. Fix basis vectors $\langle v_1\rangle=(L^\prime_1)_c$ and $\langle v_2 \rangle=(L^\prime_2)_c$ then let $V_t:=\langle tv_1+v_2 \rangle.$ 
We obtain $\lim_{t\to 0} V_t=V^\prime$ and $\lim_{t\to \infty} V_t=V_\infty$. Then we see that $\calH_{V^\prime}(E^\prime,\Phi^\prime)=(E(-c),\Phi)$ and  $\calH_{V_\infty}(E^\prime,\Phi^\prime)=(E^{\prime\prime},\Phi^{\prime\prime})$. Here $$E^{\prime\prime}=L_0^{\prime\prime}\oplus L_1^{\prime\prime}=L_0\oplus L_1(-2c)$$ and $\Phi^{\prime\prime}=\left(\begin{array}{cc}0&0\\b^{\prime\prime}&0\end{array}\right).$ For $t\neq 0,\infty$ the Hecke transforms $E^\prime_{V_t}$ are no longer direct sums but extensions \beq \label{firstext}0\to L_0(-c)\to E^\prime_{V_t}\to L_1(-c)\to 0,\eeq or \beq \label{secondext}0\to L_1(-2c) \to E^\prime_{V_t}\to L_0\to 0.\eeq The first one corresponds to  the modification of $E^\prime_1=L_0^\prime=L_0$ in $E^\prime_{V_t}$. In other words \eqref{firstext} is the modified filtration of Proposition~\ref{modifiedfilt}. This shows that $$\lim_{\lambda\to 0}\lambda\cdot \calH_{V_t}(E^\prime,\Phi^\prime) = (E(-c),\Phi).$$ In particular, as the latter is stable so is $\calH_{V_t}(E^\prime,\Phi^\prime)$.  
\source{very-stable-higgs-bundles-equivariant-multiplicity-mirror-symmetry}

The second extension \eqref{secondext} is induced from $L_1(-2c)\subset E^\prime_{V_t}$ the Hecke modification of the subbundle $L_1^\prime\cong L_1(-c)\subset E^\prime$. As $\Phi^\prime$ was trivial on $L_1^\prime$ the modification $\Phi^\prime_{V_t}$ will be trivial on $L_1(-2c)$. Thus $\Phi^\prime_{V_t}$ will be given by projection to $L_0$ in \eqref{secondext} followed by $b^{\prime\prime}:L_0\to L_1(-2c)K$. As in Proposition~\ref{downwardfiltration} this shows that $$\lim_{\lambda\to \infty} \lambda\cdot \calH_{V_t}(E^\prime,\Phi^\prime) =(E^{\prime\prime},\Phi^{\prime\prime}). $$

Thus $\calH_{V_t}(E^\prime,\Phi^\prime)$ is a $
\T$-equivariant family of stable nilpotent Higgs bundles parametrized by $t\in \P^1=\C\cup \infty$, connecting $\calH_{V_0}(E^\prime,\Phi^\prime)=(E(-c),\Phi)$ and $\calH_{V_\infty}(E^\prime,\Phi^\prime)=(E^{\prime\prime},\Phi^{\prime\prime})$.

This shows that $(E(-c),\Phi)$ and thus $(E,\Phi)$ are not very stable. 

\subsubsection{Hecke transforms of Lagrangians}
We have seen above how the different  upward flows of  type $(1,\dots,1)$ very stable Higgs bundles are related by Hecke transforms. We also showed they were Lagrangian in Proposition~\ref{lagrangian}. It is a general fact that the Hecke transform takes Lagrangians to Lagrangians, as we show next.

Simpson constructed \cite[Theorem 4.10]{simpson1} the fine moduli space $\calR^s$ of  stable rank $n$ degree $d$ Higgs bundles, framed at the point $c\in C$. The group $\GL_n$ acts on the framing with quotient  $\M^s$ so that $\calR^s\to \M^s$ is a $\PGL_n$-principal bundle, associated to the projective universal bundle $\P(\E_c)$ restricted to $c$. Let $H_k\subset \GL_n$ be the stabilizer of $\C^k<\C^n$ - a maximal parabolic subgroup of $\GL_n$. The quotient  $\GL_n/H_k$ is isomorphic to the Grassmannian of $k$-planes in $\C^n$ and we can construct the Grassmannian bundle \beq \label{grassmann} Gr_k(\E_c):=\calR^s/H_k\eeq over $\M^s$. 
\source{very-stable-higgs-bundles-equivariant-multiplicity-mirror-symmetry}

By identifying $K_c\cong \C$, \'etale locally the universal Higgs field $\bPhi_c$ restricted to $c$ gives an endomorphism $\bPhi_c\in H^0(\M^s;\End(\E_c))$. We will denote the set of $\Phi_c$-invariant $k$-dimensional subspaces of $E_c$ by $$\calH^s:=\{V< \E_{((E,\Phi),c)} | \mbox{ s.t. } 
(E,\Phi)\in \M^s, \dim(V)=k \mbox{ and } \Phi_c(V)\subset V \}\subset Gr_k(\E_c).$$ The subset $\calH^s$ has the structure of a variety, as \'etale-locally in $\M^s$ we can assume that  some $\lambda\in\C$ is never an eigenvalue of $\Phi_c$ in other words $\det(\Phi_c-\lambda)\neq 0$. That means that $\bPhi_c-\lambda$ is invertible and thus acts on $\Gr_k(\E_c)$ and  $\calH^s$ in this \'etale neighbourhood is the fixed point variety of a morphism. By construction \beq\label{proper}\pi:\calH^s\to \M^s\eeq is a proper map.
\source{very-stable-higgs-bundles-equivariant-multiplicity-mirror-symmetry}

A point of $\calH^s$ is represented by a pair $((E,\Phi),V) \in \calH^s$  of a Higgs bundle $(E,\Phi)$ and $k$-subspace $V\in Gr_k(E_{c})$ such that $\Phi_c(V)\subset V$.   We also consider ${}^{s}\calH^s\subset \calH^s$ the subset of points $((E,\Phi),V)\in \calH^s$ where $V\in Gr_k(E_{c})$ such that $\calH_{V}(E,\Phi)$ is stable. As stability is an open condition for Higgs bundles ${}^{s}\calH^s\subset \calH^s$ is an open subset. Further we define $\M^\#\subset\M$ the open subset where the spectral curve is smooth and the Higgs field $\Phi_c$ at the point $c\in C$ has distinct eigenvalues. Then we set $\calH^\#:=\pi^{-1}(\M^\#)\subset \calH^{s}$ an open subset.  The Hecke transform does not change the spectral curve, and when it is smooth the Higgs bundle is automatically stable. Thus $\calH^\#\subset {}^{s}\calH^s$. 

By abuse of notation we will still denote $\pi:\calH^\# \to \M^\#$ the restriction of \eqref{proper}. If  $\M=\M_{n,d}$ then we denote $\M^\prime:=\M_{n}^{d-k}$ so that $\calH_V(E,\Phi)\in \M^{\prime s}$ when $((E,\Phi),V)\in {}^{s}\calH^s$.  This defines $\pi^\prime:\calH^\#\to \M^{\prime\#}$.  We have the following
\begin{proposition}
\source{very-stable-higgs-bundles-equivariant-multiplicity-mirror-symmetry}
\notready\label{heckeinjective} The map
\source{very-stable-higgs-bundles-equivariant-multiplicity-mirror-symmetry}
	$(\pi^\prime,\pi):\calH^\#\to \M^{\prime \#}\times  \M^\#$ is injective. 
\end{proposition}
\begin{proof}
	When the Higgs field $\Phi_c$ has $n$ distinct eigenvalues $\lambda_i$, a Hecke transform depends on the choice of  a $k$-element subset  or equivalently a subspace $\C^k< \C^n$ spanned by the corresponding eigenspaces.  If the equation of the spectral curve is given by $\det (x-\Phi)=0$ and $(E,\Phi)$ is the direct image of the line bundle $U$ on $C_a$ then the Hecke transform is the direct image of $U(-D)$  where $D$ is the effective divisor $(\lambda_{i_1},c)+(\lambda_{i_2},c)+\dots +(\lambda_{i_k},c)$.
	
	The map will be injective if for two choices  of subspace the divisor classes $D_1,D_2$ are distinct. This will follow if $\dim H^0(C_a, {\mathcal O}(D))=1$.
	By Riemann-Roch and Serre duality this condition is equivalent to $\dim H^0( C_a, K_{C_a}(-D))=\tilde g-k$ where $\tilde g$ is the genus of $ C_a$, or alternatively if  the restriction $H^0(C_a, K_{ C_a})\rightarrow H^0(D, K_{ C_a})\cong \C^k$ is surjective. 
	
	By adjunction on the total space of $K$, which has trivial canonical bundle, we have $K_{ C_a} \cong \pi^*K^n$ and hence sections of $\pi^*K^n$ of the form $a_0+a_1x+\dots +a_{n-1}x^{n-1}$ where $a_i$ is the pullback of a section of $H^0(C,K^{n-i})$. This space has dimension $n^2(g-1)+1=\tilde g$ and so represents uniquely  every section of $K_{ C_a}$. Restricting to $D$ gives $(v_1,\dots, v_k)$ where 
	$$v_m=a_0(c)+a_1(c)\lambda_{i_m}+\dots +a_{n-1}(c)\lambda_{i_m}^{n-1}.$$
	Since the $\lambda_i$ are distinct, Lagrange interpolation provides a polynomial of degree $(n-1)$ which agrees with any $(v_1,\dots, v_k)$ and because $K^i$ on $C$ has no base points, there exist sections $a_i$ agreeing at $c$ with the coefficients of the polynomial. Hence restriction is surjective.
\end{proof} 

We can define a correspondence on ${\mathcal M}'^s\times {\mathcal M}^s$ by saying that two points $(m,m')$ are equivalent if they are represented by Higgs bundles $(E,\Phi), (E',\Phi')$ where $(E',\Phi')$ is related to $(E,\Phi)$ by a $k$-plane Hecke transformation at $c$. Then, from Proposition~\ref{heckeinjective},  
at a generic point this is locally given by the graph of a map and hence has a dense open set which is a manifold of half the dimension. 





Recall that a {\em Lagrangian correspondence} between symplectic manifolds $M_1,M_2$ is a Lagrangian submanifold  $ M_1\times M_2$ with respect to the symplectic form $p_1^*\omega_1-p_2^*\omega_2$, for example the graph of a symplectic transformation. Then points in $M_2$ corresponding to those in a Lagrangian in $M_1$ will, under transversality conditions, form a Lagrangian submanifold. The following result was already discussed in \cite[pp.185-186]{kapustin-witten}. 

\begin{theorem}
\source{very-stable-higgs-bundles-equivariant-multiplicity-mirror-symmetry}
\notready The Hecke correspondence on ${\mathcal M}'^s\times {\mathcal M}^s$ is Lagrangian at a generic point.
\end{theorem}

\begin{proof}
	We use a differential-geometric approach and assume that $\Phi_c$ has distinct eigenvalues. 
	As defined above, when we have  a rank $n$ vector bundle $E$ on the curve $C$ and a point $c\in C$ the Hecke transform $E'$ is the vector bundle defined by the kernel subsheaf of a homomorphism $h: {\mathcal O}(E)\rightarrow {\mathcal O}_c^{n-k}$ with the kernel at $c$ the vector space $V$. 
	
	
	Let $v_1,\dots, v_n$ be a local basis of sections of $E$  with $v_1,\dots, v_k$ spanning $V$ at $c$. Using a local coordinate where $c$ is $z=0$,  a local basis for $E'$ is given by $ u_1,\dots,  u_n$ where $v_i=u_i$ for $i>k$  and $v_i=zu_i$ for $1\le i\le k$. 
	We shall use these two local bases to make calculations.
	
	Varying $V\subset E_c$ gives a variation of the holomorphic structure on $E'$ as in \cite{NR}. 
	A first order deformation of $V$  in a parameter $t$  is described by local holomorphic sections of the form 
	$$ v_i+t\dot v_i=v_i+t\sum_{j=k+1}^n a_{ij}v_j$$
	for $1\le i\le k$. The variation of a local basis of the  subsheaf is then 
	$$\dot u_i=\frac{1}{z} \sum_{j=k+1}^n a_{ij}u_j.$$
	The matrix $a=a_{ij}$ in the basis $u_1,\dots, u_n$ for $E'$ gives the    section
	$$\begin{pmatrix}  0 & 0 \\
	a/z& 0\end{pmatrix}$$ over a punctured disc. 
	
	For a fixed holomorphic structure on $E$, this is   a \v Cech
	cocycle for a class in $H^1(C;\End (E'))$ defining the variation of the holomorphic structure on the Hecke transform corresponding to $a(c)\in  \Hom(V,E_c/V)$, a tangent vector to the Grassmannian.    Here however we need to vary $E$ also and it is convenient to use 
	a Dolbeault representative $\dot A\in \Omega^{0,1}(C;\End (E))$. We may make the assumption that the deformation of holomorphic structure on $E$ is trivial in a neighbourhood of $c$ so that we can take $\dot A=0$ in this neighbourhood,  extend $a$ to a $C^{\infty}$ section  supported in a neighbourhood of $c$ and holomorphic in a smaller one and take $\dot A+\bar\partial \dot a/z\in \Omega^{0,1}(C; \End (E'))$ giving the combined deformation of holomorphic structure on $E'$. 
	
	
	Let $(\dot A,\dot\Phi)$ be a first order deformation of a Higgs bundle preserving the deformation of the subsheaf.   
	A  local section of $\End (E)$ can be written in block matrix form with respect to the subspaces  spanned by $v_1,\dots, v_k$ and $v_{k+1},\dots, v_n$ respectively as
	$$\begin{pmatrix}  P & Q\\
	R & S\end{pmatrix} $$ and  in the basis $u_1,\dots, u_n$ it becomes 
	\begin{equation}
	\begin{pmatrix}  P & zQ\\
	z^{-1}R & S .\end{pmatrix}
	\label{matrix1}
\source{very-stable-higgs-bundles-equivariant-multiplicity-mirror-symmetry}
	\end{equation} 
	From this we see explicitly that  a Higgs field $\Phi\in H^0(C;\End (E)\otimes K)$  extends to $\End (E')\otimes K$ if  and only if  the entry $R$ is divisible by $z$, i.e. $\Phi$ preserves $V$ at $c$, which is the definition of $\calH^s$.
	
	
	
	
	
	To preserve the subsheaf to first order we need 
	$(\Phi+t\dot\Phi)(v+ta(v))=w+ta(w)$ modulo $t^2$ 
	for $v$ and $w$ linear combinations of  $v_1,\dots,v_k$. It follows that  $\dot\Phi+[\Phi,a]$ preserves the subspace $V$ at $c$.
	
	
	Write $$\Phi=\begin{pmatrix}  P & Q\\
	R & S\end{pmatrix}\quad \dot\Phi=\begin{pmatrix}  \dot P & \dot Q\\
	\dot R & \dot S\end{pmatrix}$$
	Since $\dot A=0$ in a neighbourhood the entries are holomorphic.  
	Then
	$$\dot\Phi+[\Phi,a]=\begin{pmatrix}  \dot P- Qa& \dot Q\\
	\dot R +Sa-aP& \dot S-aQ\end{pmatrix}.$$
	This  preserves $V$, so $\dot R +Sa-aP$ must be divisible by $z$.  Since $\Phi_c$ has distinct eigenvalues $\lambda_i$ we can take the basis vectors $v_i$ to be eigenvectors of $\Phi$ in a neighbourhood and then the condition is that 
	$\dot R_{ij}+(\lambda_i-\lambda_j)a_{ij}$ be divisible by $z$ and accordingly  $a_{ij}(0)$ is uniquely determined by $\dot \Phi$.
	
	In the basis for $E'$ we have  
	$$\dot\Phi+[\Phi,a]=\begin{pmatrix}  \dot P- Qa& \dot Q z\\
	(\dot R +Sa-aP)z^{-1}& \dot S-aQ\end{pmatrix}.$$
	which is now holomorphic in a neighbourhood of $c$.
	
	
	
	Let $(\dot A, \dot \Phi)\in \Omega^{0,1}(C;\End (E))\oplus \Omega^{1,0}(C,\End (E))$ be a Dolbeault representative of a tangent vector to ${\mathcal M}^s$. The holomorphic symplectic form $\omega$ is defined by
	$$\int_C\tr(\dot A_1\dot\Phi_2)-\tr(\dot A_2\dot\Phi_1).$$
	
	
	
	With $\dot A'=\dot A+\bar\partial  a/z$ as above we obtain
	$$ \int_C\tr(\dot A'_1\dot\Phi'_2)= \int_C\tr(\dot A_1\dot\Phi_2)+\tr (\bar\partial  a_1Q_2/z)-\tr \,\bar\partial  a_1[\Phi,a_2]/z$$
	But $Q_2$ is holomorphic on the support of $ a_1$ and divisible by $z$ so by Stokes' theorem the second term vanishes. By the Jacobi identity 
	$\tr (a_1[\Phi,a_2])-\tr (a_2[\Phi,a_1])=\tr (\Phi[a_1,a_2])$ but the off-diagonal matrices $a_1,a_2$ commute so  $\tr\, \bar\partial  a_1[\Phi,a_2]/z=\tr \,\bar\partial  a_2[\Phi,a_1]/z$.
	Then
	$$ \int_C\tr(\dot A'_1\dot\Phi'_2)= \int_C\tr(\dot A_1\dot\Phi_2)$$
	and the two symplectic structures coincide. 
\end{proof}
